
Data attribution methods that identify influential training examples are used for model debugging, data valuation, and machine learning auditing. Different methods produce different answers when applied to the same model and test example. This work investigates whether this disagreement reflects methodological deficiency or a structural phenomenon. Through systematic evaluation of four attribution methods (TRAK, TracIn, Influence Functions, FastIF) under controlled compute budgets on CIFAR-10 with a ResNet-18 model, the experiments demonstrate that attribution quality is multi-dimensional: methods face trade-offs between rank preservation and magnitude fidelity. In convex settings (logistic regression), cross-metric partial correlations exceed 0.99 at all tested compute levels, indicating tight metric coupling. In non-convex deep networks, this coupling breaks down, with R-squared from regressing metrics on approximation error norm dropping to 0.034. Different methods identify largely non-overlapping sets of influential examples, with top-50 Jaccard similarity as low as 0.0024 (less than 1\% overlap). These findings suggest that the practical question for method selection depends on which quality dimension is relevant for a given application.
